\section{The model and its separator messages}\label{sec:model}
For $n\ge1$, consider
\begin{equation}\label{eq:problem}
 v_\xi=\min\left\{x^TQx+c^Tx+(\lambda+\xi)^Tz:
 x_i(1-z_i)=0\ (i\in[n]),\quad z\in\{0,1\}^n\right\}.
\end{equation}
Here $[n]=\{1,\ldots,n\}$, $Q\in\Q^{n\times n}$ is symmetric, and
$c,\lambda\in\Q^n$. We are given rational bounds
\begin{equation}\label{eq:spectral}
 0<\mu I\preceq Q\preceq HI,\qquad |c_i|\le C.
\end{equation}
Penalties may have either sign. A support means the set $A=\{i:z_i=1\}$;
it need not equal $\{i:x_i\ne0\}$. In particular, an active coordinate
may be zero. All norms below are Euclidean unless specified otherwise.

The support graph of $Q$ has vertices $[n]$ and edges $ij$ for $i\ne j$
with $Q_{ij}\ne0$. A tree decomposition consists of a tree and vertex
sets called bags. Every graph vertex belongs to a bag, every graph edge
has both endpoints in a bag, and the bags containing any given vertex
form a connected subtree. Its width is the largest bag size minus one.
We assume a decomposition of width $w$ with $b$ bags is supplied.

For disjoint sets $I,S\subseteq[n]$ and a box radius $R\ge0$, define the
conditional message
\begin{equation}\label{eq:message}
 M_{I,S}(t)=\min_{\substack{x_I\in\R^I,\ z_I\in\{0,1\}^I\\
                     x_i(1-z_i)=0\ (i\in I)}}
 \left\{x_I^TQ_{II}x_I+(c_I+2Q_{IS}t)^Tx_I
           +\sum_{i\in I}(\lambda_i+\xi_i)z_i\right\},
 \qquad t\in[-R,R]^{|S|}.
\end{equation}
Boundary unary terms, boundary penalties, and edges within $S$ are
excluded. This convention keeps the internal matrix equal to the
positive-definite principal matrix $Q_{II}$.

For $A\subseteq I$, completing the square gives
\begin{align}
 x_A(t)&=-\tfrac12Q_{AA}^{-1}(c_A+2Q_{AS}t),\label{eq:optimizer}\\
 q_A(t)&=-\tfrac14(c_A+2Q_{AS}t)^TQ_{AA}^{-1}
                    (c_A+2Q_{AS}t)+\sum_{i\in A}\lambda_i,\label{eq:branch}\\
 F_A(t)&=q_A(t)+\sum_{i\in A}\xi_i,\label{eq:noisybranch}\\
 M_{I,S}(t)&=\min_{A\subseteq I}F_A(t).
 \label{eq:envelope}
\end{align}
The empty support contributes zero. Coordinates outside $A$ are zero.
Every $q_A$ is a full quadratic polynomial with rational coefficients, as
is $F_A$ when the perturbations are rational;
``full'' means defined on the whole parameter space, rather than restricted
to a region where it is optimal. A dictionary is any subset of the noisy branches $F_A$
whose minimum equals \eqref{eq:envelope} on the specified box.

Root the decomposition. For a child--parent edge, let $S$ be the bag
intersection and let $I$ contain the vertices in the child subtree outside
$S$. The running-intersection property implies that no graph edge joins
$I$ to a vertex outside $I\cup S$. Indeed, take a bag containing such an
edge. If it lies in the child subtree, both endpoints belong to $I\cup S$.
If it lies outside, the endpoint in $I$ occurs on both sides of the
child--parent edge, forcing it into $S$. Both cases give a contradiction.
Thus \eqref{eq:message} is its exact
subtree message. We also request $M_{[n],\varnothing}$, whose constant value
is the global optimum. These are at most $b$ messages, each with
$|S|\le w+1$.

Boundary indicators need no separate family of messages. Setting a boundary
bit to zero restricts the corresponding coordinate of $t$ to zero. Setting
that bit to one permits the same zero value, with its penalty accounted for
outside the message. Hence a single dictionary handles all boundary states
by restriction to coordinate slices.
